% Klaro -- Study Guide / Viva Prep
% Companion to main.tex (the submitted report). Not for submission -- for personal revision.
\documentclass[11pt,a4paper]{article}
\usepackage[margin=2.2cm]{geometry}
\usepackage{xcolor}
\usepackage[breakable]{tcolorbox}
\usepackage{enumitem}
\usepackage{longtable}
\usepackage{booktabs}
\usepackage{array}
\usepackage{titlesec}
\usepackage{parskip}
\usepackage{hyperref}
\hypersetup{colorlinks=true, linkcolor=blue!50!black, urlcolor=blue!50!black}

\definecolor{keyblue}{RGB}{30,39,97}
\newtcolorbox{keynums}[1][]{colback=blue!4,colframe=keyblue,boxrule=0.6pt,arc=2pt,
  left=8pt,right=8pt,top=5pt,bottom=5pt,breakable,title={Key numbers to remember},
  fonttitle=\bfseries\small,coltitle=white,colbacktitle=keyblue,#1}
\newtcolorbox{bigpicture}[1][]{colback=black!3,colframe=black!55,boxrule=0.5pt,arc=2pt,
  left=8pt,right=8pt,top=5pt,bottom=5pt,breakable,title={Big picture},
  fonttitle=\bfseries\small,coltitle=white,colbacktitle=black!55,#1}

\titleformat{\section}{\Large\bfseries\color{keyblue}}{\thesection}{0.6em}{}
\titleformat{\subsection}{\large\bfseries}{\thesubsection}{0.6em}{}
\titleformat{\subsubsection}{\bfseries}{\thesubsubsection}{0.6em}{}

\title{\textbf{Klaro --- Study Guide}\\[4pt]\large Page-by-page revision notes for the Major Project Part-1 report}
\author{Personal revision notes -- Yasharth Singh}
\date{}

\begin{document}
\maketitle
\vspace{-1em}

\begin{keynums}[title={How to use this guide}]
This file walks through the submitted report (\texttt{major.pdf}) in the same order, chapter by
chapter, page by page. Each chapter section below starts with the \textbf{report page number} so you
can flip between this guide and the actual report. Read the Glossary first -- every technical term
used anywhere in the report is defined there. The last section explains, one by one, what every
paper in the References list actually says.
\end{keynums}

\tableofcontents
\newpage

% =====================================================================
\section{Glossary -- every term used in the report}

\subsection{Abbreviations (report p. ix)}
\begin{description}[style=nextline,leftmargin=1.6cm]
  \item[4B] 4 billion parameters -- a way of stating how big an AI model is. More parameters
    generally means the model is more capable but also bigger to download and slower to run.
  \item[AI] Artificial Intelligence.
  \item[CMMI] Capability Maturity Model Integration -- a certificate that proves a software company
    follows a mature, disciplined development process. RFPs sometimes require bidders to hold one.
  \item[GB] Gigabyte -- unit of file/download size.
  \item[LoRA] Low-Rank Adaptation -- a cheap way to further train (``fine-tune'') a large AI model
    without retraining all of its weights. See Section~\ref{sec:lora} for the full idea.
  \item[PDF] Portable Document Format -- the file format RFPs come in.
  \item[RAG] Retrieval-Augmented Generation -- a technique where the AI first looks up the most
    relevant pieces of text and is only shown those, instead of the whole document. See
    Section~\ref{sec:rag}.
  \item[RFP] Request for Proposal -- the tender document Klaro analyzes.
  \item[TF-IDF] Term Frequency--Inverse Document Frequency -- a simple, classic way to score how
    important a word is to a document: a word scores high if it appears often in \emph{this}
    document but rarely across other documents.
\end{description}

\subsection{Nomenclature (report p. x)}
\begin{description}[style=nextline,leftmargin=1.6cm]
  \item[Agent] An AI helper with one job, e.g.\ ``write the summary.'' Klaro's Phase~1 has one agent
    (the Summary agent); the full plan has three (Requirements, Executive Summary, Starter Pack --
    Chapter 7).
  \item[Part (chunk)] A piece of the RFP, about 1{,}800 words (roughly four pages), that the AI reads
    in one go. Long RFPs are cut into several parts because the AI model cannot read an 84-page
    document all at once.
  \item[Fact] One short statement pulled out of the RFP, always kept together with the page number it
    came from.
  \item[Executive summary] The short, nine-heading overview of the RFP that Klaro produces -- the
    main deliverable.
  \item[Pipeline] The fixed sequence of steps every RFP goes through, start to finish (Figure 4.3:
    extract $\to$ split $\to$ analyze $\to$ combine $\to$ write $\to$ check $\to$ show).
  \item[Local model] An AI model that runs on the user's own computer, never sending data over the
    internet -- as opposed to a cloud AI service like a public chatbot.
  \item[Parameter] One of the internal numbers an AI model is made of (its learned weights). More
    parameters usually means a bigger, slower, but often more capable model.
  \item[Word piece (token)] The small unit of text an AI model actually reads and writes -- often
    part of a word rather than a whole word. Model speed is usually measured in tokens (``word
    pieces'') per second.
  \item[Backed by facts] Describes a summary point that matches a fact actually found in the RFP
    (as opposed to something invented).
  \item[Made-up information] A statement the AI produced that is \emph{not} in the source document.
    Also called \textbf{hallucination} -- see Section~\ref{sec:hallucination}.
  \item[Precision] Of everything the AI listed as a requirement, the share that were actually
    correct. High precision = few wrong answers.
  \item[Recall] Of everything that \emph{should} have been listed, the share the AI actually found.
    High recall = few missed answers.
  \item[Benchmark] A fair test that gives several AI models the exact same task so their scores can
    be compared directly.
\end{description}

\subsection{Extra terms used in the report body (not in the official lists)}
\begin{description}[style=nextline,leftmargin=2.2cm]
  \item[F1-score] A single number that balances precision and recall (their harmonic mean). Used in
    Table 6.1 because looking at precision or recall alone can be misleading -- a model that lists
    everything gets perfect recall but terrible precision, and F1 punishes that.
  \item[Chunking / chunk overlap] Splitting a long document into parts (chunks). Klaro overlaps
    neighbouring parts by 200 words so a sentence that happens to fall right on the cut boundary is
    not lost or split awkwardly (report p.\,6, Section 4.5).
  \item[Issuer detection] Klaro's way of working out \emph{who published} the RFP: it looks for the
    organization name that appears most often, with a bias toward names ending in words like
    ``Bank,'' ``Authority,'' or ``Ministry'' (report p.\,8, Section 5.4).
  \item[Grounding] A quality-check score (Figure 6.3) measuring what share of summary points are
    actually backed by an extracted fact.
  \item[Traceability] A quality-check score measuring what share of summary points carry a correct
    page number.
  \item[Coverage] A quality-check score measuring what share of the collected facts made it into the
    final summary.
  \item[Duplication] A quality-check score measuring how much repeated content is in the summary.
  \item[TextRank] A graph-based algorithm (see Section~\ref{sec:textrank}) for picking out the most
    important sentences or keywords in a text, used by Trappey et al.\ [2] and explained fully by
    Mihalcea \& Tarau [3].
  \item[Transformer / self-attention] The neural-network architecture behind essentially every
    modern AI language model, including Qwen, Gemma and Phi (see Vaswani et al.\ [5],
    Section~\ref{sec:refs}).
  \item[Fine-tuning] Continuing to train an already-trained AI model on new, task-specific examples
    so it gets better at that specific task. LoRA (below) is a cheap way to do this.
\end{description}

% =====================================================================
\newpage
\section{Chapter 1: Introduction (report p. 1)}

\begin{bigpicture}
An RFP is long (50--100 pages) and its important details are scattered across it. You cannot safely
upload a confidential RFP to a public AI chatbot. Klaro's answer: run a small AI model
\emph{entirely on the user's own laptop} and produce a nine-heading, page-referenced summary.
\end{bigpicture}

\subsection*{1.1 General}
Defines what an RFP is and lists the things someone must find in it by hand today: issuer, work
required, who can bid, deadlines, payment terms, contract length, special conditions. States the two
real test RFPs used throughout the report are 54 and 84 pages long. Explains \emph{why} cloud AI is
out (confidentiality) and states Klaro's core idea: the AI model runs locally, so the document
``never leaves it'' (Figure 1.1).

\textbf{Figure 1.1 (``Klaro in one picture'')}: shows the RFP PDF going into the Klaro desktop app,
which talks to a small AI model, both boxed inside ``your laptop.'' A separate box for ``Online AI
services'' sits outside with a dashed arrow crossed out and labelled ``never uploaded'' -- the single
clearest visual statement of Klaro's privacy design.

\subsection*{1.2 Problem Statement}
One formal sentence version of the whole project: given an RFP PDF, produce a short, nine-heading,
page-referenced summary; must run on an ordinary laptop using only a local AI model; must never send
any part of the document over the internet; must fail loudly (not silently) when something goes
wrong.

\subsection*{1.3 Objectives (memorize these four)}
\begin{enumerate}[leftmargin=1.4cm]
  \item Build the step-by-step pipeline (PDF $\to$ page-referenced summary).
  \item Pick the best small AI model by testing several on the same task.
  \item Measure summary quality and report weaknesses honestly.
  \item Package it as a desktop app for Windows and macOS with progress updates and a side-by-side
    original-document view.
\end{enumerate}

\subsection*{1.4 Scope of the Present Work}
This phase = English, text-based PDFs only, one feature (the executive summary). German is the
\emph{first planned} additional language (mid-term goal, not done yet). Scanned files, Word files,
requirement checklists, proposal drafts and searching past proposals are all future work (Chapter 7).

\begin{keynums}
54 and 84 pages (the two real test RFPs) $\cdot$ nine fixed summary headings $\cdot$ 4 objectives
$\cdot$ scope = English + text PDFs only, this phase.
\end{keynums}

% =====================================================================
\newpage
\section{Chapter 2: Background Study (report p. 2--3)}

\begin{bigpicture}
Five research papers were studied that are ``close to our problem.'' None of them do what Klaro
does: summarize a single confidential tender, on a laptop, with every point traceable to a page.
Table 2.1 on p.\,3 is the one-line summary of this whole chapter -- learn that table.
\end{bigpicture}

\subsection*{2.1 General}
Two eras of document analysis: older methods counted words and picked important sentences (extractive,
statistical); newer methods use AI language models that can read and understand long text
(abstractive, generative). Klaro is in the second camp, but borrows ideas from the first.

\subsection*{2.2 Earlier Work on Tender Documents}
\begin{description}[leftmargin=2.6cm,style=nextline]
  \item[Beason et al.\ {[1]}] Processed many US government RFPs to spot technology trends. Found:
    RFPs have no consistent layout and are full of noise (contents pages, boilerplate, legal
    text) that must be cleaned out first; and structured, exact facts are more useful to readers
    than a free-flowing paragraph summary. \textbf{Klaro takes both lessons}: it skips contents
    pages and produces fixed headings instead of free prose.
  \item[Trappey et al.\ {[2]}] Summarized over 1{,}300 power-transformer tenders using TF-IDF
    (keyword scoring) + TextRank (sentence ranking). Keeping the top 35\% of sentences preserved
    about 80\% of the important requirements and cut reading time to about 29\% of the original.
    \textbf{Klaro borrows}: summarize topic-by-topic rather than the whole document at once, and aim
    for a genuine reduction in length (this is exactly the weakness Klaro found in its own long-RFP
    results -- Section 6.4.2).
  \item[Knisely \& Pavliscsak {[4]}] Sorted research proposals into themes using
    \emph{semi-supervised} clustering: a small number of expert-labelled examples guided the
    computer while it sorted the rest, which worked much better than sorting with no human input at
    all. \textbf{Klaro connection}: Klaro's own results showed the AI sometimes files a fact under
    the wrong heading, which is why Phase 2's plan (Section 7.3.3) starts with building a small
    human-checked (``golden'') reference set.
\end{description}

\subsection*{2.4 Reading Very Long Documents}
Standard AI models can only look at a limited amount of text at once, and the computational cost
grows quickly as the text gets longer (this is the quadratic cost of self-attention -- see Vaswani
et al.\ [5] below). Longformer [6] fixes this so a model can read thousands of words in one go.
Douzon et al.\ [7] go further for business documents by also using page layout, so a value in a table
stays linked to its heading -- but this is too heavy to run on a laptop. \textbf{Klaro's own answer}:
since a small laptop model cannot read an 84-page RFP in one go anyway, Klaro deliberately cuts the
document into parts, but keeps a page number on every single piece so nothing loses its source.

\subsection*{2.5 Small AI Models and Made-Up Answers}
Language models sometimes state things that are not in the text they were given [8, 9], and they pay
less attention to the middle of very long inputs [10] (see ``Lost in the Middle,'' Section~\ref{sec:refs}).
One standard remedy is retrieval-augmented generation (RAG) [11] -- look up only the relevant
passages first, then show the model only those. \textbf{Klaro's remedy}: force the model to cite a
page for every point, then automatically check each point against the collected facts afterward. A
model can later be specialized for a task cheaply via LoRA [12] (this is exactly the Phase-2 plan).

\subsection*{2.6 Summary of Gaps -- Table 2.1 (memorize this table)}
\begin{center}
\small
\begin{tabular}{@{}p{2.6cm}p{4.3cm}p{4.3cm}@{}}
\toprule
\textbf{Work} & \textbf{What it did} & \textbf{What it did not cover / was missing}\\
\midrule
Beason et al.\ [1] & Found technology trends across many RFPs & No summary of a single RFP\\
Trappey et al.\ [2] & Picked key sentences from equipment tenders & Only copies sentences; one industry\\
Knisely \& Pavliscsak [4] & Sorted proposals with a little expert help & Research proposals, not tenders\\
Longformer [6] & Let AI read very long text & Needs powerful hardware\\
Douzon et al.\ [7] & Used page layout for business documents & Too heavy for a laptop\\
\bottomrule
\end{tabular}
\end{center}
\textbf{One-line takeaway:} none of these works on confidential tenders on an ordinary laptop, and
none produces a summary where every point is linked to its page -- that gap is exactly what Klaro
fills.

% =====================================================================
\newpage
\section{Chapter 3: Requirement Analysis (report p. 4)}

\subsection*{3.1 Who Uses Klaro}
Primary user: a bid/proposal analyst, briefing a manager, with no programming background. The team
also runs the same pipeline from the command line to test changes.

\subsection*{3.2 What the System Must Do -- Table 3.1, R1--R11}
\begin{description}[leftmargin=1.2cm,style=nextline]
  \item[R1] Accept a text-based PDF, politely reject anything else.
  \item[R2] Read text page by page, remember page numbers.
  \item[R3] Split the document into parts small enough for the AI model.
  \item[R4] Pull out facts from each part using the local AI model.
  \item[R5] Combine facts from all parts, remove repeats.
  \item[R6] Write a summary under nine headings, with page numbers for every point.
  \item[R7] Check the quality of the summary automatically.
  \item[R8] Show progress while it works; allow cancelling.
  \item[R9] Show the summary beside the original PDF; jump to any cited page.
  \item[R10] Keep a history of past analyses.
  \item[R11] Check at start-up that the AI model is installed; explain how to fix it if not.
\end{description}

\subsection*{3.3 Quality Requirements (four, non-functional)}
Privacy (no internet used at all, even fonts are local) $\cdot$ cross-platform (same app, Windows +
macOS) $\cdot$ honest about errors (names which step failed and why) $\cdot$ easy to change (model
and folders live in one settings file) $\cdot$ repeatable (every analysis saved in its own timestamped
folder).

\subsection*{3.4 Hardware and Software Requirements -- Table 3.2}
\begin{keynums}
$\geq$8\,GB RAM (tested on Apple M4, 24\,GB) $\cdot$ Windows 10/11 or macOS (tested macOS 26.5) $\cdot$
Ollama running \texttt{qwen3.5:4b} ($\approx$3.4\,GB download) $\cdot$ Python 3.11+ with pypdf,
Pydantic, FastAPI $\cdot$ desktop app: Electron + React.
\end{keynums}

% =====================================================================
\newpage
\section{Chapter 4: Detailed Design (report p. 5--7)}

\begin{bigpicture}
Two systems exist in this report: the \emph{full planned} system (Figure 4.1, three AI agents,
a workflow engine, a search index) and the \emph{Phase 1} system actually built (Figure 4.2, one
agent, a simple 7-step pipeline, no search). Table 4.1 is the honest diff between the two, with
\emph{why} each simplification was made -- this table is a very likely viva question.
\end{bigpicture}

\subsection*{4.1 The Complete Planned System -- Figure 4.1}
Desktop app (Electron) $\to$ Local server (FastAPI) $\to$ Workflow engine (Python), which runs three
agents: \textbf{Requirements agent} (lists what the bidder must do), \textbf{Summary agent} (writes
the executive summary -- this is the only one built so far), \textbf{Starter Pack agent} (drafts a
first outline of the proposal from past proposals). Everything is stored on the laptop: analysis
records, a search index of documents, extracted text/results, an activity log.

\subsection*{4.2 What We Built in Phase 1 -- Figure 4.2}
Only the Summary agent. Flow: User $\to$ Desktop app $\to$ Local server (a queue, one job at a time,
because the laptop can only run one AI model at once) $\to$ Analysis job (the 7 steps of Figure 4.3)
$\to$ talks to the local AI model via Ollama $\to$ saves everything into one Results folder per
analysis, which the server reads back to show progress/results.

\subsection*{4.3 Plan versus What Was Built -- Table 4.1 (learn the ``Why'' column)}
\begin{center}
\small
\begin{tabular}{@{}p{2.1cm}p{2.0cm}p{2.6cm}p{4.4cm}@{}}
\toprule
\textbf{Part} & \textbf{Planned} & \textbf{Built} & \textbf{Why}\\
\midrule
Desktop app & Tauri & Electron & Needed a built-in PDF viewer for the side-by-side view; one language for the whole app\\
AI helpers & Three agents & Summary agent only & Get one right and measured first\\
Workflow engine & LangGraph [19] & A simple fixed sequence of 7 steps & One agent doesn't need a complex engine\\
Document search & Search index & Not needed yet & A summary reads the whole document anyway\\
Storage & Database & One folder per analysis & Easy to inspect; nothing to set up\\
Input files & PDF and Word & PDF only & Test RFPs were PDFs\\
User feedback & Accept/correct each point & View and check only & Planned for Phase 2\\
\bottomrule
\end{tabular}
\end{center}

\subsection*{4.4 The Document Processing Pipeline -- Figure 4.3 (the seven steps, memorize the order)}
\begin{enumerate}[leftmargin=1.4cm]
  \item \textbf{Extract text} -- read the PDF page by page.
  \item \textbf{Split into parts} -- keep sections together where possible.
  \item \textbf{Analyze each part} $\leftarrow$ uses the AI model (two focused questions per part).
  \item \textbf{Combine results} -- merge facts, remove repeats.
  \item \textbf{Write the summary} $\leftarrow$ uses the AI model (writes the nine headings).
  \item \textbf{Check quality} -- nine automatic checks.
  \item \textbf{Show the report} -- in the desktop app.
\end{enumerate}
Only steps 3 and 5 touch the AI model; the rest is ordinary, deterministic program logic (fast, always
gives the same answer for the same input).

\subsection*{4.5 Handling Long Documents -- Figure 4.4}
A small model can't read an 84-page RFP at once, so Klaro works like a team splitting up a long
report: cut into $\approx$1{,}800-word parts ($\approx$4 pages each), read separately, facts noted
with page numbers, then all combined into one list that the Summary agent turns into the final
summary. \textbf{Neighbouring parts overlap by 200 words} so a sentence cut right at a part boundary
is not lost.

\subsection*{4.6 What the Summary Contains -- Table 4.2 (the nine headings, memorize in order)}
\begin{enumerate}[leftmargin=1.4cm]
  \item Overview -- what is being bought and why, 2--3 sentences.
  \item Issuing organization -- who published the RFP.
  \item Scope -- the work and deliverables required.
  \item Key requirements -- documents/forms the bidder must submit.
  \item Eligibility -- who may bid: certificates, turnover, experience.
  \item Important dates -- deadlines, bid validity, review dates.
  \item Commercial requirements -- pricing, billing, payment terms, taxes, deposits.
  \item Project duration -- contract length and extensions.
  \item Important conditions -- penalties, confidentiality, ownership of work, exit/termination rules.
\end{enumerate}

\subsection*{4.7 Keeping the Document Safe}
The app only ever connects to programs on the same laptop (its own local server + the AI model
runner). The renderer (screen part) is locked down and can only ask for three things: the server
address, a file-picker dialog, and opening a link in the browser. The embedded PDF viewer can only
show the current analysis's own RFP -- nothing else. The summary panel can only ask to jump to a page
number that actually exists in that document.

\begin{keynums}
7 pipeline steps (2 use the AI) $\cdot$ $\approx$1{,}800-word parts, 200-word overlap $\cdot$ 9 summary
headings $\cdot$ 3 planned agents vs.\ 1 built.
\end{keynums}

% =====================================================================
\newpage
\section{Chapter 5: Implementation (report p. 8--10)}

\subsection*{5.1 Tools Used -- Table 5.1}
\begin{description}[leftmargin=2.6cm,style=nextline]
  \item[Ollama {[13]}] Runs the AI model on the laptop.
  \item[Qwen3.5 4B {[14]}] The AI model itself -- reads the RFP parts, writes the summary.
  \item[Python + pypdf {[15]}] Reads text out of the PDF, page by page.
  \item[Pydantic {[16]}] Checks the AI's answers have the expected structure (schema validation).
  \item[FastAPI {[17]}] The local server between the desktop app and the analysis logic.
  \item[Electron + React {[18]}] The desktop app itself.
\end{description}

\subsection*{5.2 Reading and Splitting the PDF}
Text read from each page, stored with its page number. Contents pages near the start are skipped.
Headings are only recognized when they \emph{look like} real headings (e.g.\ ``4. SCOPE OF WORK'' or
``SECTION-D: \dots''); ordinary all-caps words are ignored, since RFPs are full of them. Each part is
then built from whole sections where possible.

\subsection*{5.3 Asking the AI Two Focused Questions}
Early tests showed asking the model for \emph{everything} in one question gave poor results (missed
items, cut-off answers). So each part is read \textbf{twice}: pass 1 asks about overview/scope/
requirements/eligibility; pass 2 asks about dates/money/duration/legal conditions. Every answer must
follow a fixed structured format with a page number per point. Klaro tolerantly recovers usable points
from a slightly malformed answer instead of discarding it. If the model fails on \textbf{more than
half} the parts, the whole analysis is marked failed rather than silently continuing with missing
data.

\subsection*{5.4 Combining Facts and Finding the Issuer}
Facts from all parts grouped by heading, repeats removed, page numbers merged. \textbf{Issuer
detection}: Klaro looks for the organization name mentioned most often (favouring names ending in
words like Bank, Authority, Ministry). \textbf{Bug found and fixed}: the first version assumed a
fixed default organization whenever it found no match, so every RFP got wrongly credited to the
issuer of the very first test document. The fixed version now answers ``Not identified'' when unsure.

\subsection*{5.5 Writing and Checking the Summary}
Combined facts given to the model with instructions to write all nine headings, keep page numbers,
and use \emph{only} the provided facts. If a heading comes back empty despite facts existing for it,
those facts are copied in directly as a safety net so nothing is silently lost (this ``safety copy''
is exactly what causes the ``too-good'' scores discussed in Section 6.4.2). Automatic checks then ask:
does every heading exist? is every point backed by a fact? does every point have a page number? how
many facts made it in? are any points repeated? how much shorter is the summary than the raw facts?
is anything under the wrong heading? These checks compare \emph{words}, not meaning -- they confirm
the summary agrees with the extracted facts, not that the extraction itself was correct.

\subsection*{5.6 The Desktop Application -- Figures 5.1--5.3}
On open, the app starts the local server in the background and waits until it responds; if the model
is missing, a setup screen explains how to install it. \textbf{Figure 5.1} (user flow): Open Klaro
$\to$ [AI ready? no $\to$ Setup screen; yes $\to$] Choose an RFP $\to$ Watch progress (can cancel)
$\to$ Read the summary $\to$ Click a page number $\to$ PDF opens at that page $\to$ Checked.
\textbf{Figure 5.2}: home screen -- drag/drop or choose a PDF, table of past analyses with pages,
chunks, duration, status. \textbf{Figure 5.3}: progress screen -- live step list with per-step timing
and a percentage bar for the current step (``Chunk 2 of 14'').

\subsection*{5.7 Testing}
\textbf{Fifteen} automatic tests cover: report-page rendering, the order of pipeline steps
(orchestration), handling of failed parts, and issuer detection. All fifteen pass. The AI steps were
tested end-to-end on the real documents used in Chapter 6. App tested by hand on macOS; a teammate
tests Windows.

\begin{keynums}
2 focused questions per part $\cdot$ fail-analysis threshold = more than half the parts failing
$\cdot$ 15 automatic tests, all passing $\cdot$ 7 automatic quality checks on the summary.
\end{keynums}

% =====================================================================
\newpage
\section{Chapter 6: Experimental Results and Analysis (report p. 11--14)}

\begin{bigpicture}
Two separate experiments happen in this chapter: (A) a \textbf{model benchmark} to pick which AI
model Klaro should use (Section 6.1, Tables 6.1--6.3), and (B) \textbf{full end-to-end runs} of the
finished app on real RFPs (Section 6.3, Table 6.4). Do not mix these two tables up in the viva --
6.1--6.3 are about \emph{choosing the model}; 6.4 is about \emph{testing the finished product}.
\end{bigpicture}

\subsection*{6.1 Choosing the AI Model}
Three candidate small models compared: \textbf{Gemma 4 (e4b)}, \textbf{Qwen3.5 (4B)},
\textbf{Phi-4-mini}. Same 20 RFP parts given to each; asked to list requirements; compared against 47
hand-marked (ground-truth) requirements.

\textbf{Table 6.1 -- Accuracy (47 hand-marked requirements):}
\begin{center}
\small
\begin{tabular}{@{}lccc@{}}
\toprule
 & \textbf{Gemma 4 (e4b)} & \textbf{Qwen3.5 (4B)} & \textbf{Phi-4-mini}\\
\midrule
Requirements found (of 47) & 33 & 23 & 7\\
Wrong/extra items listed & 22 & \textbf{8} & 9\\
Precision & 60\% & \textbf{74\%} & 44\%\\
Recall & \textbf{70\%} & 49\% & 15\%\\
F1-score & \textbf{65\%} & 59\% & 22\%\\
Mandatory vs.\ optional, correct & 85\% & \textbf{91\%} & 57\%\\
Made-up requirements & 0\% & 0\% & 6\%\\
\bottomrule
\end{tabular}
\end{center}

\textbf{Table 6.2 -- Size, speed, memory:} Gemma = 8.0B params / 9.6\,GB download / 6.7s per part
typical; Qwen3.5 = 4.7B params / 3.4\,GB / 8.1s per part; Phi-4-mini = 3.8B params / 2.5\,GB / 1.3s
per part (fastest). Peak memory: Gemma 4.4\,GB, Qwen 6.0\,GB, Phi 4.3\,GB.

\textbf{Table 6.3 -- Following the answer format:} all three models were usable 100\% of the time
after Klaro's tolerant reading, but Qwen followed the \emph{exact} requested format 0\% of the time
(it gave a bare list 30\% of the time) versus Gemma 45\% exact / Phi 35\% exact.

\subsubsection*{Why Qwen3.5 was chosen (this is the single most important paragraph in Chapter 6)}
\begin{itemize}[leftmargin=1.2cm]
  \item Phi-4-mini ruled out immediately: fastest, but found only 7/47 requirements and was the
    \emph{only} model that invented (hallucinated) a requirement.
  \item Between Gemma and Qwen: Gemma found more and was faster (higher F1, 65\% vs.\ 59\%), \textbf{but}
    it is a much bigger model -- 8B params, 9.6\,GB, almost 3$\times$ Qwen's 3.4\,GB. Every user has
    to download and store this, so a smaller model matters for a laptop app.
  \item Qwen was the \emph{more careful} model: 74\% precision vs.\ Gemma's 60\%, only 8 wrong items
    vs.\ Gemma's 22, and best at judging mandatory vs.\ optional (91\%).
  \item \textbf{Core argument}: for a summary people rely on, a \emph{wrong} point does more harm
    than a \emph{missing} one -- a missing point can be recovered by asking more focused questions
    (which is exactly why Klaro asks two question-sets per part), while a wrong point may go
    unnoticed.
  \item Qwen's weakness (not following exact format) is handled separately, by Klaro's tolerant answer
    reading -- so it isn't a real cost.
  \item Because model choice is one config setting, Gemma can be swapped in any time for comparison;
    with improved instructions its score later rose to 67\%, so it remains the main alternative.
\end{itemize}

\subsection*{6.2 First Experiments and the Pilot Run}
Before the two-question design existed, the very first attempt asked for \emph{all} requirements in
one single question per part. Result on 5 parts: 2 empty answers, 1 cut off/unusable, single answers
took 2--10 minutes. This directly caused the two-question redesign. The first full pilot run (by a
teammate, old single-question-per-part-analysis but two-question design) on the 84-page RFP: 51
points, every one backed by facts with a page number, 0 made-up info, 76\% of facts included, summary
36\% shorter than the facts. \textbf{9 points were filed under the wrong heading} (e.g.\ a CMMI
certificate under ``conditions'' instead of ``eligibility'') -- this is the concrete evidence behind
the Chapter 2 discussion of Knisely \& Pavliscsak.

\subsection*{6.3 Full Runs through the Desktop App -- Table 6.4}
Finished app run on \textbf{two RFPs across three runs}: 84-page government RFP for software
development services (R1); 54-page bank RFP for a technical service provider, run \textbf{twice} (R2,
R3 -- to test run-to-run consistency).
\begin{center}
\small
\begin{tabular}{@{}lccc@{}}
\toprule
 & \textbf{R1} & \textbf{R2} & \textbf{R3}\\
\midrule
Pages & 84 & 54 & 54\\
Parts analyzed & 17 & 14 & 14\\
Total time & 22m 51s & 25m 19s & 22m 16s\\
Points in summary & 170 & 211 & 216\\
Backed by facts / page numbers & 100\% / 100\% & 100\% / 100\% & 100\% / 100\%\\
Facts included in summary & 99\% & 97\% & 96\%\\
Shorter than the facts by & 8\% & 8\% & 9\%\\
Made-up information & none & none & none\\
Parts where AI failed & 0 & 0 & 0\\
\bottomrule
\end{tabular}
\end{center}

\subsection*{6.4 What the Results Show}
\begin{description}[leftmargin=3.4cm,style=nextline]
  \item[6.4.1 Time Taken] $\approx$95\% of total time is step 3 (AI reads each part twice,
    $\approx$1.5 min/part); writing the summary adds $\approx$1 more minute. So a 50--85 page RFP =
    22--25 minutes total.
  \item[6.4.2 Why Some Scores Look Too Good] \textbf{The most important finding in the whole report.}
    Comparing files showed every point under the seven list headings was an \emph{exact copy} of a
    collected fact -- for long RFPs the model only wrote the Overview heading itself, and the
    ``safety copy'' fallback from Section 5.5 filled in the rest. So for a long RFP, the ``summary''
    right now is really a cleaned-up list of facts, not a genuine summary -- and the ``100\% backed
    by facts'' score is high \emph{partly because of that copying}, not purely because the model is
    accurate. Likely cause: the fact list for a long RFP is too long for the amount of text the
    model reads at once. Not yet confirmed -- it's the first fix planned for Phase 2.
  \item[6.4.3 Same Document, Different Answers] R2 and R3 = identical PDF, but the AI found 236 facts
    the first time and 244 the second, with some heading counts differing by up to 10. Lesson: AI
    output is not perfectly deterministic, so improvements must be judged over \emph{several} runs,
    not one.
  \item[6.4.4 Issuer Detection] Concrete before/after: with the old code, the bank RFP (R2/R3) would
    have been wrongly credited to the government agency from R1 (the first test document). After the
    fix, it correctly returned ``India Post Payments Bank Limited,'' with the ten pages where that
    name appears.
  \item[6.4.5 Checking the Summary] Every point's page number is clickable and opens the source PDF
    at that page (Figure 6.2). The Quality tab (Figure 6.3) shows Grounding, Traceability, Coverage,
    Duplication scores, with an explicit note that these compare \emph{words}, not meaning. No
    internet connection was made during any run.
\end{description}

\begin{keynums}
Model chosen: \textbf{Qwen3.5 (4B)}, 74\% precision, 49\% recall, 59\% F1, 3.4\,GB. $\cdot$
2 RFPs / 3 runs $\cdot$ 22--25 min per run $\cdot$ 100\% points backed \& page-numbered, every run
$\cdot$ 0 made-up info, every run $\cdot$ biggest known weakness: long-RFP summaries are really just
copied facts (Section 6.4.2).
\end{keynums}

% =====================================================================
\newpage
\section{Chapter 7: Conclusion and Future Scope (report p. 15--16)}

\subsection*{7.1 Conclusion}
One-paragraph project summary: Phase 1 works end-to-end (54--84 page tender in, 22--25 minutes later
a nine-heading, page-referenced summary out, nothing sent over the internet). Model chosen by testing
three on the same task; Qwen3.5 won on fewest wrong claims and $\approx$1/3 the size of its closest
competitor. Same pipeline runs from the CLI (for testing) and inside the desktop app (for analysts);
failures are surfaced clearly. Testing exposed the long-RFP copying weakness, which sets the agenda
for Phase 2.

\subsection*{7.2 Limitations (five -- be ready to state all five)}
\begin{enumerate}[leftmargin=1.4cm]
  \item Only three RFPs tested so far -- can't yet generalize to tenders in general.
  \item No human-written reference summaries exist yet to compare against.
  \item Long-RFP summaries aren't genuinely shortened yet (Section 6.4.2).
  \item AI answers vary run to run, and 22--25 minutes is slow.
  \item Only English, text-based PDFs work; tables lose their layout.
\end{enumerate}
(The report also notes: no installer yet -- the app runs from the project folder.)

\subsection*{7.3 Future Scope}
\begin{description}[leftmargin=3.6cm,style=nextline]
  \item[7.3.1 Proposal Assistant (Figure 7.1)] Two planned modes. \textbf{Agentic workflow}: three AI
    helpers share one reading of the RFP -- Requirements agent (checklist), Executive Summary agent
    (existing), Starter Pack agent (drafts a first proposal outline with risks/questions).
    \textbf{File search workflow}: past proposals/case studies are indexed locally; the most similar
    ones to a new RFP are found and handed to the Starter Pack agent.
  \item[7.3.2 Real-Time Progress Updates] Show partial results \emph{as} each step produces them
    (facts appearing live, summary streaming as it's written) instead of only a step name.
  \item[7.3.3 Better Summaries and Better Testing] (a) Fix the long-document copying problem: read
    more at once, write one heading at a time, record when the safety copy is used. (b) Build a
    30--40 RFP reference set: a larger LLM drafts a first summary, a human corrects it, the corrected
    version becomes the ``golden'' summary used to test every future change. (c) Fine-tune Qwen3.5
    with LoRA on those golden summaries (cheap fine-tuning, not full retraining) and compare against
    the current model -- this directly targets the 6.4.2 copying problem. (d) Repeat the model
    benchmark, keeping Gemma 4 as the main alternative.
  \item[7.3.4 App Improvements] One-click installers (Win + macOS) with guided model setup; support
    for Word files, scanned PDFs, tables; \textbf{multilingual support starting with German}
    (read German RFPs, reply in German or English -- first added language, mid-term goal); let users
    accept/edit/flag each point and use that feedback to measure quality.
\end{description}

\begin{keynums}
5 limitations $\cdot$ Proposal Assistant = 2 workflows (agentic + file search), 3 agents total
$\cdot$ Phase-2 plan = fix copying + 30--40 RFP golden set + LoRA fine-tune + repeat benchmark $\cdot$
German = first planned language.
\end{keynums}

% =====================================================================
\newpage
\section{References -- what every paper actually says (report p. 17)}
\label{sec:refs}

\begin{bigpicture}
References [1]--[12] are genuine research papers; [13]--[19] are tool/library documentation pages
(Ollama, Qwen3.5, pypdf, Pydantic, FastAPI, Electron, LangGraph), not papers -- listed here for
completeness since they're in the References list too.
\end{bigpicture}

\subsection*{[1] Beason, Hinton, Salamah \& Salsman -- ``Automated analysis of RFPs using NLP for the
technology domain'' (SMU Data Science Review, 2021)}
Processed a large set of real US Department of Defense and Homeland Security RFPs with NLP, looking
for technology trends across the whole collection (not summarizing any one RFP). Their two big
findings, both directly reused by Klaro: RFPs follow no consistent layout and contain a lot of noise
(contents pages, boilerplate, legal text) that must be stripped out first; and readers get more value
from exact, structured facts than from a free-flowing prose summary.

\subsection*{[2] Trappey et al.\ -- ``Intelligent RFQ summarization using NLP, text mining, and
machine learning techniques'' (Journal of Global Information Management, 2022)}
Summarized 1{,}300+ power-transformer procurement tenders (RFQs) using a two-stage pipeline: TF-IDF
to find important keywords, then TextRank to rank and select sentences by topic. Keeping only the
top 35\% of ranked sentences preserved about 80\% of the important requirements while cutting reading
time to roughly 29\% of the original. Klaro borrows the ``summarize per topic, not all at once'' idea
and the goal of a genuine length reduction (a goal Klaro has not yet met for long RFPs -- see
Section 6.4.2).

\subsection*{[3] Mihalcea \& Tarau -- ``TextRank: Bringing order into text'' (EMNLP, 2004)}
\label{sec:textrank}
The foundational graph-based ranking algorithm cited by [2]. Builds a graph where sentences (or
words) are nodes, connected by edges weighted by similarity or co-occurrence; then runs a
PageRank-style random-walk algorithm over that graph to score how ``central''/important each node is.
The highest-scoring sentences (or words) become the extractive summary (or keyword list). It needs no
training data and no deep model -- purely statistical/graph-based, which is why it is cheap enough to
run on 1,300+ documents in [2].

\subsection*{[4] Knisely \& Pavliscsak -- ``Research proposal content extraction using NLP and
semi-supervised clustering'' (Scientometrics, 2023)}
Sorted a large set of research grant proposals into themes using semi-supervised clustering: a small
number of proposals were labelled by human experts, and those labels guided an algorithm in sorting
the rest automatically. This worked substantially better than clustering with no human guidance at
all. Klaro connects this to its own finding that the AI sometimes misfiles a fact under the wrong
heading (Section 6.2) -- both point toward a little targeted human correction going a long way, which
is exactly why Phase 2 starts with a small human-checked reference set rather than pure automation.

\subsection*{[5] Vaswani et al.\ -- ``Attention is all you need'' (NeurIPS, 2017)}
The paper that introduced the \textbf{Transformer} architecture -- the design behind essentially every
modern AI language model, including Gemma, Qwen and Phi, the three models Klaro benchmarks. Its core
idea is \textbf{self-attention}: instead of reading text strictly left-to-right (like older recurrent
models), every word can directly look at every other word in the input to decide what matters,
enabling parallel training and much better handling of relationships between distant words. The
catch, relevant to Klaro's whole chunking design: the computational cost of self-attention grows
\emph{quadratically} with input length, so a standard Transformer can only look at a limited amount of
text at once -- exactly the constraint cited in report Section 2.4 that forces Klaro to split long
RFPs into parts.

\subsection*{[6] Beltagy, Peters \& Cohan -- ``Longformer: The long-document transformer'' (arXiv,
2020)}
Redesigns self-attention so its cost grows \emph{linearly} (not quadratically) with input length, by
combining local attention (each word attends mainly to nearby words) with a small amount of global
attention (a few special tokens that attend to everything). This lets a Transformer read thousands of
words in a single pass. Report's honest caveat: this still needs more powerful hardware than an
ordinary laptop, which is why Klaro chunks instead of adopting a Longformer-style model.

\subsection*{[7] Douzon, Duffner, Garcia \& Espinas -- ``Long-range transformer architectures for
document understanding'' (ICDAR 2023 Workshops)}
Extends the long-context idea specifically for business documents, by also feeding the model the
document's page layout (positions of headings, tables, columns) so that, for example, a value inside
a table stays correctly linked to its column/row heading rather than being read as flat text. More
capable for structured documents, but heavier to run -- report's stated reason it's ``too heavy for
a laptop.''

\subsection*{[8] Ji et al.\ -- ``Survey of hallucination in natural language generation'' (ACM
Computing Surveys, 2023)}
\label{sec:hallucination}
A broad survey paper cataloguing the different ways language-generation models produce
\textbf{hallucinations} -- statements not supported by (or contradicting) the source input -- across
many NLP tasks, along with the known causes (e.g.\ training-data noise, the model's own generation
process) and mitigation strategies proposed in the literature. Cited as the general evidence behind
Klaro's Section 2.5 statement that ``language models sometimes state things that are not in the text
they were given.''

\subsection*{[9] Maynez, Narayan, Bohnet \& McDonald -- ``On faithfulness and factuality in
abstractive summarization'' (ACL, 2020)}
An empirical study specifically on \emph{summarization} models (directly relevant to Klaro): shows
that abstractive summarizers frequently generate content that is not supported by, or directly
contradicts, the source document, and proposes ways to actually measure this ``faithfulness'' rather
than just measuring fluency. This is the direct research justification for why Klaro's automatic
quality checks (Section 5.5, Figure 6.3) exist at all -- checking that every summary point is
``grounded'' in an extracted fact.

\subsection*{[10] Liu et al.\ -- ``Lost in the middle: How language models use long contexts''
(TACL, 2024)}
Shows that language models are noticeably better at using information placed at the very
\emph{beginning} or \emph{end} of a long input than information buried in the \emph{middle} -- a
U-shaped performance curve as position in the context varies. Cited in Klaro Section 2.5 as a second
reason (besides hallucination) that simply throwing a whole long RFP at a model in one go is risky --
reinforcing the case for Klaro's chunk-by-chunk reading instead of one giant prompt.

\subsection*{[11] Lewis et al.\ -- ``Retrieval-augmented generation for knowledge-intensive NLP
tasks'' (NeurIPS, 2020)}
\label{sec:rag}
Introduced \textbf{RAG}: pair a \emph{retriever} (which searches a large corpus and fetches the most
relevant passages for a given query) with a \emph{generator} (a language model that only has to read
those retrieved passages, not the whole corpus). This directly reduces hallucination, since the model
is only ever shown text that is actually relevant, and its answers can be traced back to specific
retrieved passages. Klaro's own design (cite a page for every point, then check each point against the
collected facts) is described in the report as its own answer to the same underlying problem RAG
solves, without needing a full retrieval index for a single-document task.

\subsection*{[12] Hu et al.\ -- ``LoRA: Low-rank adaptation of large language models'' (ICLR, 2022)}
\label{sec:lora}
Introduces \textbf{LoRA (Low-Rank Adaptation)}: instead of updating all of a large pretrained model's
weights during fine-tuning (expensive, and requires storing a full new copy of the model per task),
LoRA freezes the original weights entirely and injects small, trainable \emph{low-rank} matrices into
specific layers (typically the attention projections). Only these small matrices are trained, which
drastically cuts the number of trainable parameters and the compute/memory needed to fine-tune,
while keeping inference speed almost unchanged. This is exactly the technique named in Klaro's
Phase-2 plan (Section 7.3.3) for fine-tuning Qwen3.5 on the 30--40-RFP golden-summary reference set.

\subsection*{[13]--[19] Tool documentation (not research papers)}
\begin{description}[leftmargin=2.4cm,style=nextline]
  \item[{[13]} Ollama] Official docs for the tool that runs AI models locally on the laptop.
  \item[{[14]} qwen3.5 model library entry] Ollama's page for the specific model Klaro uses.
  \item[{[15]} pypdf documentation] The Python library used to read text out of PDFs.
  \item[{[16]} Pydantic documentation] The Python library used to validate the AI's structured
    answers.
  \item[{[17]} FastAPI documentation] The Python web-framework used for Klaro's local server.
  \item[{[18]} Electron security documentation] Cited specifically for the security practices
    (context isolation, no direct Node access, etc.) behind Section 4.7.
  \item[{[19]} LangGraph documentation] The workflow-engine library that was \emph{planned} for the
    full multi-agent system (Figure 4.1) but not used in Phase 1, per Table 4.1.
\end{description}

% =====================================================================
\newpage
\section{Cheat Sheet -- every number in one place}

\begin{longtable}{@{}p{7.2cm}p{7.2cm}@{}}
\toprule
\textbf{Fact} & \textbf{Number}\\
\midrule
\endhead
Test RFP lengths & 54 and 84 pages\\
Summary headings & 9 (fixed, always in the same order)\\
Objectives & 4\\
Limitations listed & 5\\
Papers reviewed in Background Study & 5 close papers, 19 references total\\
Functional requirements & R1--R11 (11 total)\\
Pipeline steps & 7 (steps 3 and 5 use the AI model)\\
Chunk size / overlap & $\approx$1{,}800 words per part / 200-word overlap\\
Model benchmark test set & 20 RFP parts, 47 hand-marked requirements\\
Gemma 4 (e4b) & 33/47 found, 60\% precision, 70\% recall, 65\% F1, 8.0B params, 9.6\,GB\\
Qwen3.5 (4B) -- \textbf{chosen} & 23/47 found, 74\% precision, 49\% recall, 59\% F1, 4.7B params, 3.4\,GB\\
Phi-4-mini -- ruled out & 7/47 found, 44\% precision, 15\% recall, 22\% F1, only model to hallucinate (6\%)\\
Gemma's later improved score & 67\% (with better instructions) -- remains main alternative\\
Pilot run (84-page, old single-question) & 51 points, 76\% facts included, 36\% shorter, 9 misfiled\\
Full run R1 (84 pg, govt) & 22m 51s, 170 points, 99\% facts included, 8\% shorter\\
Full run R2 (54 pg, bank, 1st) & 25m 19s, 211 points, 97\% facts included, 8\% shorter, 236 facts found\\
Full run R3 (54 pg, bank, 2nd) & 22m 16s, 216 points, 96\% facts included, 9\% shorter, 244 facts found\\
Made-up info across all 3 full runs & 0 (none)\\
Points backed by facts / page numbers, all 3 runs & 100\% / 100\%\\
Time breakdown & $\approx$95\% of time = AI reading (step 3), $\approx$1.5 min/part\\
Automatic tests & 15 (all pass)\\
Automatic quality checks per summary & 7 questions (headings exist? backed? page numbers? coverage?
  duplicates? shorter-by-how-much? right heading?)\\
Quality-tab scores (Figure 6.3 example) & Grounding 100\%, Traceability 100\%, Coverage 97\%,
  Duplication 3.8\%\\
Phase-2 fine-tuning dataset plan & 30--40 RFPs, human-corrected golden summaries, LoRA fine-tuning\\
Minimum RAM & 8\,GB (tested on 24\,GB Apple M4)\\
\bottomrule
\end{longtable}

\end{document}
